\section{Grain Synthesis}
\label{sec:grain}

\subsection{Grain Is Not Noise, and It Is Not a Texture}

Two failure modes dominate grain in existing products. The first overlays a
scanned grain plate, which is spatially fixed, resolution-dependent, uncoupled
to image content, and repeats visibly. The second adds Gaussian noise of
constant variance, which is uncoupled to density and has a flat spectrum,
producing something that reads as sensor noise rather than emulsion grain.

Real grain has four properties that must all be reproduced:

\begin{enumerate}
\item It is \emph{signal-dependent}: variance depends on how much of the image
was developed at that point, peaking in the mid-scale.
\item It is \emph{spatially correlated} at a scale set by the emulsion's grain
size, with a characteristic non-flat Wiener spectrum.
\item It is \emph{physically sized}: a grain is a few micrometres at the film
plane, so its appearance at a given viewing size depends on the format and the
enlargement, not on the pixel count.
\item It is \emph{chromatically structured}: in color film, the three dye
clouds form in separate layers and are essentially statistically independent,
producing colored grain; in monochrome film there is one layer, producing
neutral grain.
\end{enumerate}

\EM{} derives all four from a single generative model.

\subsection{The Filtered Poisson Model}

Model developed grains as a spatial Poisson point process. Let $\{\mathbf{p}_i\}$
be grain centers with intensity $\lambda$ (grains per unit area), and let each
grain contribute a density profile $a_i\, g(\mathbf{p} - \mathbf{p}_i)$, where
$g$ is the normalized dye cloud profile and $a_i$ is a random amplitude drawn
from the grain size distribution. The density field is the shot-noise process

\begin{equation}
D_g(\mathbf{p}) \;=\; \sum_i a_i\, g\!\left(\mathbf{p} - \mathbf{p}_i\right).
\label{eq:shotnoise}
\end{equation}

Campbell's theorem \cite{campbell1909} gives the first two moments and the power
spectrum in closed form:

\begin{align}
\mathbb{E}[D_g] &= \lambda\,\mathbb{E}[a] \int g(\mathbf{p})\, d\mathbf{p},
\label{eq:campbell-mean}\\
\mathrm{Var}[D_g] &= \lambda\,\mathbb{E}[a^2] \int g^2(\mathbf{p})\, d\mathbf{p},
\label{eq:campbell-var}\\
W(\mathbf{f}) &= \lambda\,\mathbb{E}[a^2]\, \left|G(\mathbf{f})\right|^2,
\label{eq:wiener}
\end{align}

where $G = \mathcal{F}\{g\}$ and $W$ is the Wiener (noise power) spectrum. This
is the entire theoretical content of the grain engine, and it is the reason the
model can be validated against published Wiener spectra (AC-3) rather than
tuned by eye.

\subsection{Selwyn Granularity Recovered}

Measuring density through a scanning aperture of area $A$ averages the field,
which for $A$ large compared to the grain gives

\begin{equation}
\sigma_A^2 = \frac{1}{A}\,\lambda\,\mathbb{E}[a^2] \int g^2\, d\mathbf{p},
\end{equation}

so $\sigma_A \sqrt{2A}$ is a constant of the emulsion. This is
\emph{Selwyn's law} \cite{selwyn1935}, and the constant

\begin{equation}
G_S = \sigma_A\sqrt{2A} = \sqrt{2\lambda\,\mathbb{E}[a^2]\textstyle\int g^2}
\label{eq:selwyn}
\end{equation}

is the Selwyn granularity, the quantity manufacturers publish (typically as RMS
granularity measured through a $48\,\mu$m aperture, scaled by $1000$). Because
$G_S$ is published for most stocks, \eqref{eq:selwyn} lets us fit
$\lambda\mathbb{E}[a^2]$ directly from datasheet values rather than guessing.

\subsection{Density Dependence}

A pure Poisson model predicts $\sigma^2 \propto \lambda$, and since mean density
is also $\propto \lambda$ by Nutting's relation \cite{nutting1913}

\begin{equation}
D = 0.4343\,\frac{n\,\bar{a}}{A},
\end{equation}

it predicts $\sigma \propto \sqrt{D}$, monotonically increasing. Real emulsions
do not behave this way at high density: as development approaches completion the
available crystals are exhausted, grain overlap saturates, and granularity
falls again.

The correct model replaces the Poisson count with a \emph{binomial} count: there
are $N$ crystals available, and each develops with probability $p$ determined by
exposure. Then $\mathrm{Var}[n] = Np(1-p)$ and

\begin{equation}
\boxed{\;\sigma_D^2(D) \;=\; \sigma_{\mathrm{ref}}^2 \cdot \frac{p(1-p)}{p_{\mathrm{ref}}(1-p_{\mathrm{ref}})}\;}
\qquad p = \frac{D - \Dmin}{\Delta D}
\label{eq:grainvar}
\end{equation}

with $p_{\mathrm{ref}} = 0.5$ so the reference is the peak. This recovers the
observed inverted-U granularity--density relationship with no free parameters
beyond the peak magnitude.

A shape exponent pair $(\nu_1, \nu_2)$ generalizes \eqref{eq:grainvar} to
$\sigma_D^2 \propto p^{\nu_1}(1-p)^{\nu_2}$ for stocks whose measured curve is
asymmetric; monochrome stocks typically fit $(\nu_1, \nu_2) = (1.0, 0.45)$,
which rises to a peak at $p \approx 0.69$ and falls only modestly thereafter,
matching black-and-white behavior.

\begin{figure}[htbp]
\centering
\begin{tikzpicture}
\begin{axis}[
  width=\columnwidth, height=4.2cm,
  xlabel={fractional density $p = (D-\Dmin)/\Delta D$},
  ylabel={$\sigma_D$ (normalized)},
  xmin=0, xmax=1, ymin=0, ymax=1.15,
  grid=both, grid style={draw=emLine!20},
  legend style={font=\scriptsize, at={(0.5,0.05)}, anchor=south, draw=emLine!40, legend columns=2},
  label style={font=\scriptsize}, tick label style={font=\scriptsize},
  samples=200,
]
\addplot[emAcc, very thick, domain=0.001:0.999] {sqrt(4*x*(1-x))};
\addlegendentry{color, $(1,1)$}
\addplot[emBlue, very thick, domain=0.001:0.999] {(x^1.0 * (1-x)^0.45)/ (0.69^1.0 * 0.31^0.45)};
\addlegendentry{mono, $(1,0.45)$}
\end{axis}
\end{tikzpicture}
\caption{Granularity versus density from \eqref{eq:grainvar}. Grain vanishes in clear film and in fully developed maximum density, and peaks in the mid-scale --- which is why grain is most visible in skies and skin, not in shadows.}
\label{fig:grainvar}
\end{figure}

\subsection{The Grain Kernel and Clustering}

The dye cloud profile $g$ is modelled as a mixture of two Gaussians, capturing
both the individual grain and the tendency of grains to cluster:

\begin{equation}
g(\mathbf{p}) = (1-\chi)\, \mathcal{N}(\mathbf{p};\, \sigma_1) + \chi\, \mathcal{N}(\mathbf{p};\, \sigma_2),
\qquad \sigma_2 = \zeta \sigma_1,
\label{eq:grainkernel}
\end{equation}

with $\chi \in [0, 0.4]$ the clustering fraction and $\zeta \approx 3$. Since the
Fourier transform of a Gaussian is a Gaussian, the Wiener spectrum follows in
closed form from \eqref{eq:wiener}:

\begin{equation}
W(f) = \lambda \mathbb{E}[a^2]\left[(1-\chi)e^{-2\pi^2\sigma_1^2 f^2} + \chi e^{-2\pi^2\sigma_2^2 f^2}\right]^2 .
\label{eq:wienerclosed}
\end{equation}

Equation \eqref{eq:wienerclosed} is what AC-3 compares against measured
data. Fitting a stock is therefore a two-parameter least-squares problem in
$(\sigma_1, \chi)$ against its published Wiener spectrum, with $\lambda
\mathbb{E}[a^2]$ pinned by the published RMS granularity through
\eqref{eq:selwyn}. This is a genuine fit to measurement, satisfying the
"defaults are measurements" principle of Section~\ref{sec:intro-thesis}.

\subsection{Format Scaling}
\label{sec:formatscale}

All grain geometry is specified in micrometres at the film plane. The
conversion to render pixels is

\begin{equation}
\sigma^{\mathrm{px}} = \sigma^{\mu\mathrm{m}} \cdot \frac{W_{\mathrm{px}}}{1000\, W_{\mathrm{mm}}},
\label{eq:formatscale}
\end{equation}

where $W_{\mathrm{px}}$ is the render width in pixels and $W_{\mathrm{mm}}$ the
simulated frame width in millimetres. A \emph{format} selector supplies
$W_{\mathrm{mm}}$: $36$\,mm for 135, $56$\,mm for 645/6$\times$6, $102$\,mm for
4$\times$5, $24.9$\,mm for Super\,35, $12.5$\,mm for Super\,16, $10.3$\,mm for
8\,mm. Selecting Super\,16 makes the grain of the same emulsion roughly three
times larger relative to the frame, which is exactly the real relationship and
requires no separate "grain size" authoring per format.

The same equation governs $\sigma_1, \sigma_2$ of the interlayer model
\eqref{eq:twoscale} and the halation scattering length \eqref{eq:halpsf}, so
format selection coherently rescales every spatial phenomenon at once. This is
a significant structural advantage of specifying everything in physical units.

\subsection{Resolution Independence and Determinism}

NFR-08 requires that identical parameters produce identical output, and the
preview/export equivalence of NFR-07 requires that the grain \emph{structure} be
the same field sampled at two resolutions, not two different random fields.

Both are achieved by defining the noise on a lattice in \emph{film-plane
coordinates} rather than pixel coordinates, the same device that makes
stochastic grain models resolution-independent \cite{newson2017}. For a render pixel at position
$\mathbf{q}$ pixels, the film coordinate is $\mathbf{m} = \mathbf{q} \cdot
1000\,W_{\mathrm{mm}} / W_{\mathrm{px}}$ in micrometres. The white noise field is
evaluated by hashing the integer lattice cell of $\mathbf{m}$ at the grain
lattice pitch, so:

\begin{itemize}
\item The same film position yields the same random value at every render
resolution.
\item Preview at $1/4$ resolution samples the same field more coarsely, which
--- through \eqref{eq:formatscale} --- correctly reduces the per-pixel variance
according to Selwyn's law. The preview grain therefore looks like the export
grain viewed at preview size, which is what WYSIWYG actually means here.
\item Changing the frame seed changes the hash salt, producing an independent
realization, so "re-roll grain" is a one-integer operation.
\end{itemize}

\subsection{Chromatic Structure}

In color film the three dye clouds form in physically separate layers from
independent crystal populations, so the three noise fields are close to
statistically independent, producing visibly colored grain. In monochrome film
there is a single silver image, so the three output channels receive the
identical field, producing neutral grain.

\EM{} interpolates between these through a correlation matrix
$\mathbf{R}$ with unit diagonal and off-diagonal $\varphi$:

\begin{equation}
\mathbf{n} = \mathbf{L}\, \mathbf{z}, \qquad \mathbf{L}\mathbf{L}^\top = \mathbf{R},
\label{eq:graincorr}
\end{equation}

where $\mathbf{z}$ is a vector of three independent standard normal fields and
$\mathbf{L}$ is the Cholesky factor, computed once on the host. Setting
$\varphi = 0$ gives fully chromatic grain, $\varphi = 1$ gives neutral grain,
and the shipping defaults are $\varphi \approx 0.15$ for color negative,
$\varphi \approx 0.35$ for transparency (whose layers are more closely coupled),
and $\varphi = 1$ for monochrome.

Because $\mathbf{R}$ is a $3\times3$ correlation matrix, its Cholesky factor is
lower triangular with six nonzero entries, so \eqref{eq:graincorr} costs six
multiply-adds.

\subsection{Application to Density}

The grain field is added in the density domain, scaled by the density-dependent
standard deviation:

\begin{equation}
\boxed{\;\hat{D}_c(\mathbf{p}) \;=\; D_c(\mathbf{p}) \;+\; g \cdot \sigma_{D,c}\!\left(D_c(\mathbf{p})\right) \cdot n_c(\mathbf{p})\;}
\label{eq:grainapply}
\end{equation}

with $g$ the user's grain amount ($0$ to $2$, default $1$; $1$ means "the
physically correct amount for this stock"). Setting $g = 0$ produces a
grainless render, which is not physical but is a legitimate user choice; the UI
labels the $g = 1$ position on the slider explicitly.

Note that grain is added to the \emph{negative} density, which is then printed
through \eqref{eq:printcurve}. The print stock's gamma of $2.8$ therefore
amplifies negative grain by that factor in the print, and the print's toe and
shoulder \emph{suppress} grain in print highlights and print blacks. Both are
correct, both are free, and neither is achievable by adding grain to the final
image.

\subsection{Generation Algorithm}

Two synthesis modes are required, selected by the expected number of grains per
render pixel $\Lambda = \lambda\,\Delta A_{\mathrm{px}}$:

\begin{itemize}
\item \textbf{Gaussian mode} ($\Lambda \geq 10$). The central limit theorem
applies; generate a white Gaussian field and convolve with $g$. This is the
common case at full resolution.
\item \textbf{Poisson mode} ($\Lambda < 10$). The Gaussian approximation fails
and produces grain that is too smooth and insufficiently "clumpy"; individual
grains must be realized. This occurs for large formats rendered small, and for
heavily enlarged small formats. The shader draws $n \sim \mathrm{Poisson}(\Lambda)$
per lattice cell using Knuth's method \cite{knuth1997} (acceptable since
$\Lambda$ is small) and splats.
\end{itemize}

\begin{lstlisting}[language=Metal, caption={Deterministic hash and grain field generation (Gaussian mode).}, label={lst:grain}]
// PCG-style integer hash: deterministic, well-distributed, cheap.
inline uint pcg_hash(uint v) {
    uint state = v * 747796405u + 2891336453u;
    uint word  = ((state >> ((state >> 28u) + 4u)) ^ state) * 277803737u;
    return (word >> 22u) ^ word;
}
inline float hash01(uint3 c, uint salt) {
    uint h = pcg_hash(c.x ^ pcg_hash(c.y ^ pcg_hash(c.z ^ salt)));
    return float(h) * 2.3283064365386963e-10f;   // 1/2^32
}
// Box-Muller from two hashed uniforms; returns standard normal.
inline float gauss(uint3 c, uint salt) {
    float u1 = max(hash01(c, salt),        1e-7f);
    float u2 =     hash01(c, salt ^ 0x9E3779B9u);
    return sqrt(-2.0f * log(u1)) * cos(6.2831853f * u2);
}

kernel void grain_white_field(
    texture2d<float, access::write> outTex [[texture(0)]],
    constant GrainParams&           P      [[buffer(0)]],
    uint2 gid [[thread_position_in_grid]])
{
    if (gid.x >= outTex.get_width() || gid.y >= outTex.get_height()) return;

    // Map pixel -> film-plane micrometres -> grain lattice cell.
    // P.umPerPixel = 1000 * frameWidthMM / renderWidthPX
    float2 um   = float2(gid) * P.umPerPixel;
    int2   cell = int2(floor(um / P.latticePitchUM));

    float3 z;
    for (uint c = 0; c < 3; ++c)
        z[c] = gauss(uint3(uint2(cell), c), P.seed);

    // Correlate the three records via the Cholesky factor of R.
    float3 n;
    n.x = P.L00 * z.x;
    n.y = P.L10 * z.x + P.L11 * z.y;
    n.z = P.L20 * z.x + P.L21 * z.y + P.L22 * z.z;

    outTex.write(float4(n, 1.0f), gid);
}
\end{lstlisting}

The white field is then convolved with the two-Gaussian kernel of
\eqref{eq:grainkernel} using the separable recursive filter of
Section~\ref{sec:recursive-blur}, and combined with the density image by
\eqref{eq:grainapply}. Total cost is two separable blur passes plus two
pointwise passes over a single-precision RGBA texture.

\subsection{A Note on Downsampling}

If a full-resolution render is exported and then downsampled by an external
tool, the grain will be averaged and its variance reduced by the sampling
area --- correctly, per Selwyn. If instead the user exports at a reduced
resolution, \EM{} renders at that resolution and \eqref{eq:formatscale} gives
smaller pixel-space grain, again correctly. The two paths agree to within
the difference between the resampling filter and the grain kernel, which is
below visual threshold. This consistency is verified by test V-11 in
Section~\ref{sec:validation}.
